\section*{2+2=4}

We establish the arithmetic identity $2 + 2 = 4$ from first principles using the
Peano axioms, where the natural numbers are built from a constant $0$ and a
successor function $S$.

\subsection*{Definitions}

Following Peano, we define:
\[
  1 \coloneqq S(0), \quad
  2 \coloneqq S(1) = S(S(0)), \quad
  3 \coloneqq S(2) = S(S(S(0))), \quad
  4 \coloneqq S(3) = S(S(S(S(0)))).
\]

Addition is defined recursively by:
\[
  m + 0 = m, \qquad m + S(n) = S(m + n).
\]

\subsection*{Derivation}

\begin{align*}
  2 + 2
    &= 2 + S(1)      && \text{(definition of } 2\text{)} \\
    &= S(2 + 1)      && \text{(recursive case of addition)} \\
    &= S(2 + S(0))   && \text{(definition of } 1\text{)} \\
    &= S(S(2 + 0))   && \text{(recursive case)} \\
    &= S(S(2))       && \text{(base case: } m + 0 = m\text{)} \\
    &= S(3)          && \text{(definition of } 3\text{)} \\
    &= 4.            && \text{(definition of } 4\text{)}
\end{align*}

\subsection*{Remark}

This result is elementary but serves as a sanity check for the formal system:
any model of the Peano axioms must satisfy $2 + 2 = 4$.  In particular, no
non-standard model can violate it, since the derivation above uses only the
axioms and the explicit definitions of the numerals.
